% ============================================================================
% Dok. 264 -- Ferrotoroidizität und Quaternärspeicher: Alle Formeln
%             aus FFGFT-Perspektive reinterpretiert
%
% Referenzarbeit: Qureshi et al., Nature Communications 17:4033 (2026)
% DOI: 10.1038/s41467-026-70767-8
%
% Status: Reduktionsschicht -- strukturelle Analogien und Brückenargumente.
% Wo FFGFT einen natürlichen Rahmen für eine Formel bietet, wird dies
% vermerkt; spekulative Verbindungen sind explizit gekennzeichnet.
% ============================================================================

\section*{Vorbemerkung: Warum diese Arbeit für FFGFT bedeutsam ist}

Qureshi et al.\ (2026, \emph{Nature Communications} 17:4033) präsentieren
die erste experimentelle Demonstration eines quaternären (4-Zustands-)
nichtflüchtigen Speichers in einem rein antiferromagnetischen Einkristall
LiNi$_{0.8}$Fe$_{0.2}$PO$_4$, mittels ferrotoroidischer Domänenkontrolle
durch senkrechte elektrische und magnetische Felder. Die vier unterscheid-
baren Domänenzustände werden durch eine spontane Spinrotation unterhalb der
Phasenübergangstemperatur $T_4 = 21$\,K erzeugt und mit Sphärischer
Neutronenpolarimetrie (SNP) verifiziert.

Aus FFGFT-Sicht ist diese Arbeit aus drei unabhängigen Gründen bemerkenswert:
\begin{enumerate}
  \item Der \emph{Torus} ist nicht nur das fundamentale geometrische Objekt
        von FFGFT (T$^4$, Dok.~001/002), sondern erscheint direkt als
        physikalischer Ursprung der 4-Zustands-Struktur im Experiment:
        vier Domänenzustände = zwei toroidale Richtungen $\times$ zwei
        Orientierungszustände, eine $2\times2$-Zerlegung, die die
        Moden-Klassifikation auf einem kompakten Torus widerspiegelt.
  \item Der \emph{Polarisationsmatrix-Formalismus} (Blume-Maleev-Gleichungen)
        ist eine Rotationsmatrix-Darstellung der Zustandsentwicklung --
        strukturell identisch mit den unitären Operatoren auf dem T$^4$-
        Modenraum in FFGFT (Dok.~203, Rekursion $\mathcal{R}$).
  \item Die \emph{ferrotoroidische Ordnung} bricht gleichzeitig Raum- und
        Zeitumkehrsymmetrie -- genau die Symmetriestruktur der FFGFT-
        Windungsmoden, die sowohl eine räumliche Windungszahl als auch eine
        zeitliche Orientierung tragen.
\end{enumerate}

\noindent\textbf{Status dieses Dokuments.} Jede Formel der Referenzarbeit
wird reproduziert und reinterpretiert. Die FFGFT-Interpretationen sind
Reduktionsschicht, sofern nicht ausdrücklich als Brücke oder Kernableitung
ausgewiesen. Die Ergebnisse der Referenzarbeit werden nicht bestritten;
die Reinterpretation fügt eine geometrische Lesart hinzu, keine physikalische
Korrektur.

% ============================================================================
\section{Toroidisierungsformel (Gl.~1 der Referenzarbeit)}
% ============================================================================

\subsection*{Die Formel}

\begin{equation}
  \mathbf{T} = 2S\,\frac{\varepsilon_a}{\Omega}\cos\varphi \cdot \mathbf{b}
  \times \mathbf{c}
  \tag{Ref.~Gl.~1}
\end{equation}

\noindent wobei $S$ der Spinbetrag ist, $\varepsilon_a/\Omega$ ein
geometrischer Vorfaktor (Einheitszellenvolumen $\Omega$, Spinverschiebung
$\varepsilon_a$), $\varphi$ der Rotationswinkel der magnetischen Momente in
der $a$--$b$-Ebene, und $\mathbf{b}\times\mathbf{c}$ die toroidale Richtung
entlang der kristallografischen $c$-Achse definiert.

\subsection*{FFGFT-Interpretation}

In FFGFT ist der Torus T$^4$ das fundamentale geometrische Objekt mit vier
kompakten Dimensionen. Die Toroidisierung $\mathbf{T}$ ist die makroskopische
Projektion eines \emph{Windungsflusses}:
\begin{itemize}
  \item $S$ entspricht der Windungsquantenzahl $n$ eines Torusmodus
        (Dok.~001, Grundrelation T$\tilde{\cdot}m=1$): der Spinbetrag misst,
        wie oft der Modus um die kompakte Richtung windet.
  \item $\varepsilon_a/\Omega$ ist die \emph{Packungsdichte} der Windung
        in der Einheitszelle -- direkt analog zum Packungsdefizit
        $\xi = 4/30000$ in FFGFT, das misst, um wieviel das 3D-Raumgitter
        von einer glatten Mannigfaltigkeit abweicht (Dok.~009).
  \item $\cos\varphi$ ist die \emph{Projektion} der Windung auf die beob-
        achtbare Richtung. In FFGFT entspricht dies der Dekompaktifizierungs-
        projektion: die volle T$^4$-Windung ist einheitenlos (Schicht~1),
        und $\cos\varphi$ ist der Anteil, der in der SI-projizierten
        (Schicht~2) Beschreibung beobachtbar wird.
  \item $\mathbf{b}\times\mathbf{c}$ definiert die toroidale Achse -- in
        FFGFT ist dies die Projektionsrichtung vom 4D-Kompaktraum auf den
        3D-beobachtbaren Unterraum. Die Kreuzproduktstruktur ist dieselbe
        wie in der T$^4$-Volumenform, eingeschränkt auf einen 2-Zyklus.
\end{itemize}

\noindent\textbf{Schlüsselbeobachtung:} $\mathbf{T}$ ist unabhängig vom
Rotationswinkel $\varphi$ im Sinne, dass seine \emph{Richtung} (Vorzeichen)
sich bei Variation von $\varphi$ nicht ändert -- nur sein Betrag ändert sich.
Dies ist strukturell identisch mit dem FFGFT-Nichtabschluss-Theorem
(Dok.~193): die topologische Ladung (Windungszahl) ist invariant unter
stetigen Deformationen des Modus innerhalb einer Windungsklasse.

% ============================================================================
\section{Polarisationsmatrix (Gl.~2 und Gl.~3)}
% ============================================================================

\subsection*{Die Formeln}

\begin{equation}
  \mathbf{P}_f = \mathbf{P}\,\mathbf{P}_i
  \tag{Ref.~Gl.~2}
\end{equation}

\begin{equation}
  \mathcal{P} = \frac{1}{I}
  \begin{pmatrix}
    N^2 - M_\perp^2 & J_{nz} & J_{ny} \\
    J_{nz} & N^2 + M_{\perp y}^2 - M_{\perp z}^2 & R_{yz} \\
    J_{ny} & R_{yz} & N^2 - M_{\perp y}^2 + M_{\perp z}^2
  \end{pmatrix}
  \tag{Ref.~Gl.~3}
\end{equation}

\noindent wobei $N$ der Kernstrukturformfaktor ist, $M_\perp$ der magnetische
Wechselwirkungsvektor, $J_{ni} = 2\,\mathrm{Im}(NM_{\perp i}^*)$ die
Kern-Magnet-Interferenzterme, und $R_{yz} = 2\,\mathrm{Re}(M_{\perp y}M_{\perp z}^*)$.

\subsection*{FFGFT-Interpretation}

Die Transformation $\mathbf{P}_f = \mathbf{P}\,\mathbf{P}_i$ ist eine
\emph{Rotationsmatrix, die auf den Neutronenspinzustand wirkt}. In FFGFT
ist dies das direkte Analogon der Rekursion $\mathcal{R}$ (Dok.~203), die
auf den T$^4$-Modenvektor wirkt: der Zustandsvektor transformiert unter
einer unitären Rotation, die durch die geometrische Struktur des Streuobjekts
bestimmt ist.

Die Matrix $\mathcal{P}$ hat vier verschiedene Formen für die vier Domänen --
jede Domäne entspricht einer \emph{anderen Windungskonfiguration} auf dem
Torus, und jede Windungskonfiguration erzeugt eine andere Rotation des
Modevektors. Dies ist strukturell identisch mit den vier Windungsklassen
in FFGFT: zwei toroidale Richtungen (Zeitumkehrpaare, $\pm$ Windung)
$\times$ zwei Orientierungszustände (Raumspiegelungspaare), was eine
$\mathbb{Z}_2 \times \mathbb{Z}_2$-Gruppenstruktur ergibt.

Die Kern-Magnet-Interferenzterme $J_{ni}$ sind von Null verschieden, weil der
magnetische Strukturformfaktor $M_\perp$ in der antiferromagnetischen Phase
\emph{rein imaginär} ist (die durch Inversion verwandten Momente sind
antiparallel). In FFGFT entsprechen imaginäre Strukturformfaktoren Moden, die
auf gegenüberliegenden Seiten des Torus \emph{gegenphasig} sind -- genau die
antiferromagnetische (Kopf-an-Schwanz-)Konfiguration des ferrotoroidischen
Grundzustands.

Das Diagonalelement $N^2 - M_\perp^2$ misst das Ungleichgewicht zwischen
Kern-(Skalar-)Streuung und Magnet-(Vektor-)Streuung -- in FFGFT ist dies das
Ungleichgewicht zwischen dem Skalarfeld $E$ (Lagrangedichte
$\mathcal{L}=\xi(\partial E)^2$, Dok.~010) und den Vektorfeldkomponenten,
das die Energiedichte $T_{00} = \xi[(\partial_t E)^2 + (\nabla E)^2]$
bestimmt.

% ============================================================================
\section{Allgemeine Polarisationsmatrix mit Strahlenimperfektionen (Gl.~9)}
% ============================================================================

\subsection*{Die Formel}

\begin{equation}
  \mathcal{P} =
  \begin{pmatrix}
    \frac{p_{fx}}{p_{ix}}\frac{N^2-M_\perp^2-J_{yz}}{I_x} &
    \frac{p_{fx}}{p_{iy}}\frac{J_{nz}-J_{yz}}{I_y} &
    \frac{p_{fx}}{p_{iz}}\frac{J_{ny}-J_{yz}}{I_z} \\[6pt]
    \frac{p_{fy}}{p_{ix}}\frac{J_{nz}+R_{ny}}{I_x} &
    \frac{p_{fy}}{p_{iy}}\frac{N^2+M_{\perp y}^2-M_{\perp z}^2+R_{ny}}{I_y} &
    \frac{p_{fy}}{p_{iz}}\frac{R_{yz}+R_{ny}}{I_z} \\[6pt]
    \frac{p_{fz}}{p_{ix}}\frac{J_{ny}+R_{nz}}{I_x} &
    \frac{p_{fz}}{p_{iy}}\frac{R_{yz}+R_{nz}}{I_y} &
    \frac{p_{fz}}{p_{iz}}\frac{N^2-M_{\perp y}^2+M_{\perp z}^2+R_{nz}}{I_z}
  \end{pmatrix}
  \tag{Ref.~Gl.~9}
\end{equation}

\subsection*{FFGFT-Interpretation}

Die Faktoren $p_{fx}/p_{ix}$ sind \emph{Projektionseffizienzen} -- sie messen,
wie genau der Messapparat den Polarisationszustand auflöst. In FFGFT
entsprechen diese den \emph{Dekompaktifizierungsfaktoren}: dem Verhältnis
zwischen der beobachtbaren (SI-projizierten, Schicht~2) Größe und der
zugrundeliegenden geometrischen (einheitenlosen, Schicht~1) Größe.
Perfekte Messung ($p_{fi} = 1$) entspricht dem Schicht-1-Grenzfall;
reale Messung ($p_{fi} < 1$) der Schicht-2-Projektion mit ihrer inhärenten
Unvollständigkeit.

Die reellen Terme $R_{ni} = 2\,\mathrm{Re}(NM_{\perp i}^*)$ und die
imaginären Terme $J_{ni}$ bilden genau die Zerlegung in
\emph{gerade} (kosinus-, zeitspiegelsymmetrische) und \emph{ungerade}
(sinus-, zeitspiegelbrechende) Windungsmodenkomponenten (Dok.~203).

% ============================================================================
\section{Dzyaloshinskii--Moriya-Hamiltonoperator (Gl.~4--6)}
% ============================================================================

\subsection*{Die Formeln}

\begin{align}
  \mathcal{H}_1^{\mathrm{DM}} &= \mathbf{D}_{12}\cdot(\mathbf{S}_1\times\mathbf{S}_2)
    + \mathbf{D}_{34}\cdot(\mathbf{S}_3\times\mathbf{S}_4)
  \tag{Ref.~Gl.~4} \\[6pt]
  \mathcal{H}_{1c}^{\mathrm{DM}} &= D_{12}^c
    (S_1^a S_2^b - S_1^b S_2^a - S_3^a S_4^b + S_3^b S_4^a)
  \tag{Ref.~Gl.~5} \\[6pt]
  \mathcal{H}_{1b}^{\mathrm{DM}} &= -D_{12}^b
    (S_1^a S_2^c - S_1^c S_2^a - S_3^a S_4^c + S_3^c S_4^a)
  \tag{Ref.~Gl.~6}
\end{align}

\noindent mit $\mathbf{D}_{34} = -\mathbf{D}_{12}$ (Antisymmetrie unter
Inversion), und den Komponenten $D_{12}^b, D_{12}^c$ als Projektionen des
Dzyaloshinskii-Moriya-Vektors auf die kristallografischen $b$- und $c$-Achsen.

\subsection*{FFGFT-Interpretation}

Die Dzyaloshinskii-Moriya-Wechselwirkung ist ein
\emph{antisymmetrischer Austausch}: $\mathbf{D}\cdot(\mathbf{S}_1\times\mathbf{S}_2)$
ist ungerade unter Vertauschung der Spins 1 und 2. In FFGFT entsprechen
antisymmetrische Wechselwirkungen den \emph{nicht-kommutativen} Produkten von
Windungsmoden -- die Kreuzproduktstruktur $S_i\times S_j$ ist genau die
Struktur, die das Nichtabschluss-Theorem (Dok.~193) als topologisch
eingeschränktes Produkt von Modenlabeln verschiedener Windungsklassen
identifiziert.

Die Relation $\mathbf{D}_{34} = -\mathbf{D}_{12}$ (aufgrund der
Inversionssymmetrie) entspricht der FFGFT-Symmetrie, dass Moden, die durch
räumliche Inversion verwandt sind, entgegengesetzte Windungszahlen haben:
wenn Modus $n$ im Uhrzeigersinn windet, windet der inversionsverknüpfte
Modus $-n$ gegen den Uhrzeigersinn.

Die Komponenten $D^c$ und $D^b$ (mit $|D^b| < |D^c|$ aus der Geometrie)
spiegeln wider, dass der DM-Vektor nicht entlang einer Kristallachse
ausgerichtet ist -- er liegt unter $35°$ zur $c$-Achse in der $b$--$c$-Ebene.
In FFGFT entspricht dies der \emph{anisotropen $\xi$-Kopplung}: der
geometrische Parameter $\xi$ koppelt in verschiedenen kompakten Richtungen
unterschiedlich, mit der dominanten Kopplung in Richtung des größten
Windungsmodus (hier die ferrotoroidische $c$-Achsenrichtung).

% ============================================================================
\section{Polarisations-Messwert (Gl.~7)}
% ============================================================================

\subsection*{Die Formel}

\begin{equation}
  P_{fi} = \frac{n_{fi}^{++} - n_{fi}^{+-}}{n_{fi}^{++} + n_{fi}^{+-}}
  \tag{Ref.~Gl.~7}
\end{equation}

\subsection*{FFGFT-Interpretation}

Dies ist der quantenmechanische Erwartungswert des Polarisationsoperators,
normiert durch die Gesamtzählung. In FFGFT entspricht dies dem Verhältnis
zweier Modenpopulationen: Moden im ,,+``-Windungszustand gegenüber Moden im
,,-``-Windungszustand nach der Streuung (d.\,h. nach einmaliger Wirkung
des Rekursionsoperators $\mathcal{R}$).

Die Normierung durch $n^{++} + n^{+-}$ (Gesamtzählung) ist die FFGFT-
\emph{Dekompaktifizierung}: die absolute Modenanzahl (eine Schicht-1-Größe)
ist nicht direkt beobachtbar; nur das Verhältnis (eine dimensionslose,
Schicht-1-kompatible Größe) ist es. Dies ist dieselbe Struktur wie die
FFGFT-$T_{\mathrm{CMB}}$-Relation: nur das Verhältnis
$T_{\mathrm{CMB}}/E_\xi = (16/9)\xi^2$ ist eine reine geometrische Vorhersage;
der absolute Kelvin-Wert erfordert das SI-Einbettung.

% ============================================================================
\section{Endpolarisation mit erzeugten Komponenten (Gl.~8)}
% ============================================================================

\subsection*{Die Formel}

\begin{equation}
  \mathbf{P}_f = \mathcal{P}'\,\mathbf{P}_i + \mathcal{P}''
  \tag{Ref.~Gl.~8}
\end{equation}

\noindent wobei $\mathcal{P}'$ der rein rotatorische Anteil der
Polarisationstransformation ist und $\mathcal{P}''$ die erzeugte oder
vernichtete Polarisation (von Null verschieden, wenn der Streuprozess selbst
Polarisation erzeugt).

\subsection*{FFGFT-Interpretation}

Die Zerlegung in einen \emph{Rotations-}anteil $\mathcal{P}'$ und einen
\emph{Erzeugungs-}anteil $\mathcal{P}''$ entspricht der FFGFT-Unterscheidung
zwischen zwei Typen von Modentransformationen:
\begin{itemize}
  \item $\mathcal{P}'\,\mathbf{P}_i$: \emph{kohärente} Transformation --
        der Eingangszustand wird ohne Energieaustausch rotiert. Dies
        entspricht der FFGFT-Rekursion $\mathcal{R}$ (Dok.~203), die auf
        einen bestehenden Modus wirkt: die Windungszahl bleibt erhalten,
        nur die Phase (Orientierung) ändert sich.
  \item $\mathcal{P}''$: \emph{spontane} Polarisation -- neue Polarisation
        wird durch den Streuprozess selbst erzeugt. Dies entspricht der
        \emph{Modenerzeugung} in FFGFT: ein neuer Modus wird an einem
        T$^4$-Knoten generiert (Dok.~204, Fraktale Randbedingung).
\end{itemize}

\noindent In der spezifischen magnetischen Struktur von
LiNi$_{0.8}$Fe$_{0.2}$PO$_4$ gilt $\mathcal{P}''=0$, weil die magnetischen
Strukturformfaktoren rein imaginär sind. In FFGFT entspricht dies dem Fall,
dass die T$^4$-Geometrie \emph{statisch} ist: keine neuen Moden werden
spontan erzeugt, alle Transformationen sind rein rotatorisch (unitär). Dies
ist der statische T$^4$-Grundzustand -- genau das kosmologische FFGFT-Bild
(statisches Universum, keine spontane Modenerzeugung = keine Expansion).

% ============================================================================
\section{Die Vier-Domänen-Struktur: $2\times2$ auf T$^4$}
% ============================================================================

Das zentrale experimentelle Ergebnis der Arbeit sind vier unterscheidbare
magnetische Domänen, kontrolliert durch zwei unabhängige Stellhebel:
\begin{enumerate}
  \item \textbf{Toroidale Domänen} (2 Zustände): ausgewählt durch
        $\mathbf{E}\times\mathbf{H}$ (parallel oder antiparallel zu $c$).
  \item \textbf{Orientierungsdomänen} (2 Zustände): ausgewählt durch $H_c$
        (eine kleine $c$-Komponente des Magnetfeldes, über die DM-Energie).
\end{enumerate}

\noindent Die vier kombinierten Zustände ($a_4', a_4'', b_4', b_4''$) sind
verknüpft durch:
\begin{itemize}
  \item Zeitumkehr $1'$: $a_4 \leftrightarrow b_4$ (toroidaler Flip).
  \item Zweizählige Schraubenachse $2_1\|y$: $'\leftrightarrow''$
        (Orientierungsflip).
\end{itemize}

\subsection*{FFGFT-Interpretation}

Diese $\mathbb{Z}_2\times\mathbb{Z}_2$-Gruppenstruktur ist die direkte
experimentelle Realisierung der T$^4$-Windungsmoden-Klassifikation in FFGFT:
\begin{itemize}
  \item Das \emph{toroidale} $\mathbb{Z}_2$: entspricht den zwei möglichen
        Orientierungen einer geschlossenen Windungsschleife auf dem T$^4$
        (im/gegen Uhrzeigersinn = $\pm\mathbf{T}$). Dies ist der
        fundamentale ,,Zeitpfeil``-Freiheitsgrad in FFGFT: die intrinsische
        Zeit $\tilde{T}$ kann in zwei Richtungen auf dem kompakten Torus
        zeigen (Dok.~230, Bell-Korrelationen als T$^4$-Geometrie).
  \item Das \emph{Orientierungs-}$\mathbb{Z}_2$: entspricht den zwei
        möglichen räumlichen Orientierungen der Windungsebene im 4D-Torus.
        In FFGFT ist dies die räumliche Parität des Modus -- ob die Windung
        in der $+c$- oder $-c$-räumlichen Richtung verläuft.
\end{itemize}

\noindent\textbf{WBE-Verbindung:} Die Vier-Domänen-Struktur hat eine effektive
Dimensionalität $D_{\mathrm{eff}} = 2$ (zwei binäre Freiheitsgrade). Der
WBE-Exponent für diese effektive Dimension ist
$\alpha = D_{\mathrm{eff}}/(D_{\mathrm{eff}}+1) = 2/3$. Dies ist weder
$w=-3/4$ (Onurs 3D-Netzwerk) noch $w=-1$ ($\Lambda$CDM), sondern ein
Zwischenwert für die 2D-Natur der toroidalen Domänenkontrollfläche. Dies
deutet darauf hin, dass verschiedene effektive Dimensionen auf verschiedenen
Skalen gelten -- konsistent mit der FFGFT-Sicht, dass $\xi$ eine
fraktionale Abweichung von ganzzahliger Dimension auf allen Skalen codiert
(Dok.~242, ,,Die Skalenleiter``).

% ============================================================================
\section{Quaternärspeicher und FFGFT-Informationsstruktur}
% ============================================================================

Die Arbeit zeigt, dass 8 quaternäre Einheiten $4^8 = 65536$ Kombinationen
ergeben, verglichen mit $2^8 = 256$ für binär -- ein Faktor 256 Zunahme.
In FFGFT entspricht dies der fundamentalen Beobachtung, dass der T$^4$-Torus
\emph{mehr Information pro Modus} trägt als ein Punkt oder ein Kreis: jeder
Modus ist durch vier Windungszahlen (eine pro kompakter Dimension) gekennzeichnet,
nicht nur durch eine (binär).

Die \emph{nichtflüchtige} Natur des Speichers -- die Domänenpopulationen sind
stabil bei Null-Feld, nachgewiesen durch Messung bei $E=H=0$ -- entspricht in
FFGFT der \emph{topologischen Stabilität} der Windungsmoden: ein Modus mit
von Null verschiedener Windungszahl kann nicht auf Null-Windung zerfallen,
ohne eine Energiebarriere zu überwinden (Dok.~204, Fraktale Randbedingung).
Deshalb ist der Toroidspeicher nichtflüchtig: die topologische Ladung ist
erhalten, sofern kein konjugiertes Feld aufgebracht wird, das die
Anisotropiebarriere übersteigt.

% ============================================================================
\section{Bezug zum FFGFT-Parameter $\xi$}
% ============================================================================

Die Geometrie der Arbeit umfasst drei wesentliche Längenskalen:
\begin{itemize}
  \item Die Einheitszelldimension ($\sim 10$\,\AA): setzt die Skala der
        einzelnen Windungsschleife.
  \item Die Domänengröße ($\sim$mm): die makroskopische Kohärenzlänge der
        toroidalen Ordnung.
  \item Die Phasenübergangstemperaturen $T_4=21$\,K und $T_2=25$\,K:
        die Energieskalen für die zwei Ebenen der Domänenstruktur.
\end{itemize}

\noindent\textbf{Energieverhältnis:} Das Verhältnis der beiden Phasenübergangs-
temperaturen beträgt $T_4/T_2 = 21/25 = 0{,}840$. In FFGFT werden Verhältnisse
von Energieskalen durch Potenzen von $\xi$ kontrolliert. Eine geometrische
Herleitung dieses konkreten Verhältnisses aus der T$^4$-Geometrie ist eine
offene Frage (Reduktionsschicht).

% ============================================================================
\section{Zusammenfassung: Was das Experiment über die T$^4$-Struktur zeigt}
% ============================================================================

\begin{itemize}
  \item \textbf{Toroidale Ordnung ist real und kontrollierbar:} Das Experiment
        liefert den ersten eindeutigen mikroskopischen Nachweis von
        Vier-Zustands-Toroid-Domänenkontrolle in einem antiferromagnetischen
        Massivkristall. Dies bestätigt, dass die T$^4$-Topologie -- die Basis
        von FFGFT -- nicht nur ein mathematisches Konstrukt ist, sondern sich
        als physikalische Domänenstruktur in realen Materialien manifestiert.
  \item \textbf{Die $\mathbb{Z}_2\times\mathbb{Z}_2$-Struktur ist exakt:}
        Die vier Domänen sind präzise durch Zeitumkehr und eine räumliche
        Schraubenoperation verknüpft -- dieselbe Gruppe, die T$^4$-Windungs-
        moden in FFGFT klassifiziert.
  \item \textbf{Rein imaginäre Strukturformfaktoren = statischer T$^4$:}
        Dass $M_\perp$ rein imaginär ist (antiferromagnetische Konfiguration)
        bedeutet $\mathcal{P}''=0$: keine spontane Polarisation wird erzeugt.
        Dies ist der FFGFT-statische Grundzustand.
  \item \textbf{Nichtflüchtigkeit = topologische Stabilität:} Der Speicher
        ist nichtflüchtig, weil die Windungszahl topologisch geschützt ist.
        In FFGFT ist dies derselbe Mechanismus, der die Leptonmassen-Hierarchie
        stabilisiert: die Windungszahlen $p_e=3/2$, $p_\mu=2$, $p_\tau=5/2$
        sind topologisch geschützt (Dok.~190, P2).
\end{itemize}

\noindent\textbf{Status der Interpretationen.}
Abschnitte~1--5 (einzelne Formeln) sind \emph{Reduktionsschicht}: strukturelle
Analogien zwischen dem Formalismus der Referenzarbeit und FFGFT, ohne neue
quantitative Vorhersage. Abschnitte~6 (4-Domänen = T$^4$-Moden) und~7
(Nichtflüchtigkeit = topologische Stabilität) sind \emph{Brücken}: die
Entsprechung ist algebraisch exakt ($\mathbb{Z}_2\times\mathbb{Z}_2$-
Gruppenstruktur, Erhaltung der topologischen Ladung), aus beiden Rahmenwerken
unabhängig abgeleitet. Abschnitt~8 ($\xi$-kontrollierte Skalenhierarchie)
ist \emph{Reduktionsschicht}: die Struktur ist plausibel, aber die genaue
$\xi$-Potenz ist nicht bestimmt.

% ============================================================================
\section{Fraktale Wegverlängerung vs.\ Tired Light -- eine notwendige Abgrenzung}
% ============================================================================

Die FFGFT-Erklärung der Rotverschiebung als \emph{fraktale Wegverlängerung}
im statischen Universum (Dok.~190, P14) muss von der historischen
\emph{Tired-Light}-Hypothese (Zwicky 1929) abgegrenzt werden, die in der
Astrophysik aus drei Gründen abgelehnt wird. FFGFT vermeidet alle drei:

\begin{enumerate}
  \item \textbf{Keine Unschärfe ferner Quellen.} Tired Light erfordert
        Photonenstreuung an intervenierender Materie, was Wellenfronten
        aufweitet und ferne Galaxien unscharf erscheinen ließe. Die fraktale
        Wegverlängerung beinhaltet keine Streuung: das Photon folgt der
        topologischen Geometrie des statischen Vakuums (der in
        LiNi$_{0.8}$Fe$_{0.2}$PO$_4$ nachgewiesenen ,,Head-to-Tail``-
        Windungsstruktur) und erhält die volle Wellenfrontkohärenz.
  \item \textbf{Konsistente geometrische Zeitdilatation $(1+z)$.} Tired Light
        sagt keine Zeitdilatation von Supernova-Lichtkurven vorher, was
        observationell ausgeschlossen ist. In FFGFT trägt jedes Signal auf
        einem fraktal verlängerten Weg automatisch dasselbe Wellenlängen-
        verhältnis $\lambda_b/\lambda_e$ in Wellenlänge wie in Pulsdauer --
        beide skalieren geometrisch durch die fraktale Wegverlängerung, nicht
        durch Energieverlust. Das Verhältnis bleibt erhalten, und die
        beobachtete Supernovae-Zeitdilatation folgt zwangsläufig.
        (Das Symbol $z = \lambda_b/\lambda_e - 1$ ist dabei nur eine
        Bezeichnung für das gemessene Verhältnis, keine FFGFT-Primärgröße,
        Dok.~190, P14.)
  \item \textbf{Erhalt des CMB-Schwarzkörperspektrums.} Photonenstreuung
        würde die Spektralform verzerren. Fraktale Wegverlängerung skaliert
        alle Frequenzen gleichmäßig ($\Delta\nu/\nu = \mathrm{const}$) und
        erhält die perfekte Schwarzkörperform des CMB-Spektrums.
\end{enumerate}

\noindent Die im Papier nachgewiesene ferrotoroidische Vakuumstruktur liefert
den physikalischen Mechanismus für Punkt~(1): die ,,Head-to-Tail``-Konfiguration
(rein imaginäres $M_\perp$, daher $\mathcal{P}''=0$, keine spontane Polarisation)
sichert kohärente, streuungsfreie Ausbreitung.

% ============================================================================
\section{Topologische Struktur auf drei Skalen}
% ============================================================================

\begin{table}[htbp]
\centering
\caption{Dieselbe $\mathbb{Z}_2\times\mathbb{Z}_2$-topologische Struktur
erscheint auf drei verschiedenen Skalen.}
\label{tab:skalen}
\begin{tabular}{llll}
\toprule
\textbf{Skala} & \textbf{System} & \textbf{Topologie} & \textbf{Zustände} \\
\midrule
Atomar ($10^{-10}$\,m) &
  LiNi$_{0.8}$Fe$_{0.2}$PO$_4$ &
  Toroidale Spin-Ringe &
  4 Domänen (SNP) \\
Sub-Planck ($10^{-38}$\,m) &
  FFGFT-Vakuum (T$^4$) &
  Toroidale Windungsmoden &
  4 Windungsklassen \\
Kosmologisch (Mpc) &
  Beobachtbares Universum &
  Fraktale Wegverlängerung &
  $H_0 = 66{,}2$\,km/s/Mpc (kein $\Lambda$) \\
\bottomrule
\end{tabular}
\end{table}

\noindent Die kosmologische Zeile verwendet $H_0$ aus $\xi$ (Dok.~263/026)
als geometrische Kenngröße, keinen Zustandsgleichungsparameter $w$: FFGFT
hat kein Dunkle-Energie-Feld und keine Expansion, daher gibt es kein
zuzuweisendes $w$. Die Skaleninvarianz der topologischen Struktur ist
strukturell, nicht numerisch.

% ============================================================================
\section{Korrespondenztabelle: Papier $\leftrightarrow$ FFGFT}
% ============================================================================

\begin{table}[htbp]
\centering
\caption{Verifizierte Entsprechungen zwischen Qureshi et al.\ (2026) und FFGFT.
Nur Einträge mit expliziter Ableitung oder strukturellem Argument in beiden
Rahmenwerken. Spekulative oder numerisch nicht gestützte Gleichsetzungen
sind ausgeschlossen.}
\label{tab:korrespondenz}
\begin{tabular}{ll}
\toprule
\textbf{Nature-Papier (Qureshi et al.\ 2026)} & \textbf{FFGFT (T$^4$-Geometrie)} \\
\midrule
Head-to-Tail Spin-Ringe & Geschlossene Windungsschleifen auf T$^4$ \\
Toroidaler Moment $\mathbf{T}\parallel\hat{\mathbf{c}}$ &
  Windungsfluss in kompakter $c$-Richtung \\
Spontaner Spin-Tilt ($\varphi$) &
  Dekompaktifizierungsprojektion ($\cos\varphi$, Schicht~2) \\
4 diskrete Domänen &
  $\mathbb{Z}_2\times\mathbb{Z}_2$ Windungsmodenklassen \\
Zeitumkehr $1'$: $a\leftrightarrow b$ &
  Zeitliche Windungsrichtung ($\pm\tilde{T}$) \\
Schraubenachse $2_1\|y$: $'\leftrightarrow''$ &
  Räumliche Windungsparität ($\pm c$) \\
Polarisationsmatrix $\mathcal{P}$ (Rotation) &
  Rekursion $\mathcal{R}$ auf T$^4$-Modenvektor (Dok.~203) \\
$\mathcal{P}''=0$ (rein imaginäres $M_\perp$) &
  Statischer T$^4$-Grundzustand (keine Modenerzeugung) \\
Nichtflüchtiger Speicher &
  Erhalt der topologischen Ladung (Dok.~204) \\
Messwert $P_{fi}$ (dimensionslos) &
  Schicht-1-Verhältnis (keine SI-Einheit erforderlich) \\
\bottomrule
\end{tabular}
\end{table}

\noindent\textbf{Referenz.} N.~Qureshi et al., \textit{Toroidizität als Weg
zu nichtflüchtigem Quaternärspeicher in Antiferromagneten},
\textit{Nature Communications} \textbf{17}, 4033 (2026).
DOI:~10.1038/s41467-026-70767-8.
